\documentclass{beamer}
\usepackage{amsfonts}
\usepackage{graphicx}
\graphicspath{ {images/} }
\usepackage{caption}
\usepackage{subcaption}
%Information to be included in the title page:
\title{Groups, graphs, and morphisms}
\author{Zayn Patel \\ \texttt{zaynpatelwhs@gmail.com}}
\date{May 21, 2026}
\setbeamertemplate{navigation symbols}{}

\begin{document}

\frame{\titlepage}

\begin{frame}
\frametitle{Goal}
The goal of this talk is to see a relationship between \textbf{group} theory and \textbf{graph} theory. \bigskip

To do this we will introduce Cayley DiGraphs of Groups which will allow us to visualize a group as a directed graph. \bigskip

This visualization can help us understand properties of a group (e.g. the relations) but we will discover that Cayley DiGraphs give us many valuable theorems about groups we like (e.g. $\mathbb{Z}_{m} \oplus \mathbb{Z}_{n})$.
\end{frame}

\begin{frame}
\frametitle{Topics covered}
To accomplish our goal we will cover the following: \bigskip
    \begin{itemize}
        \item Generators and relations
        \item Cayley DiGraphs of Groups
        \item Hamiltonicity of DiGraphs
        \item Graph theory meets groups theory
        \item Generators and relations rigor
    \end{itemize}
\end{frame}

\begin{frame}
\frametitle{Generators and relations}
    \begin{itemize}
        \item<1-> Generators and relations are special elements we can combine, via the group operation, to get any other element in the group.
        \item<2-> Groups with a single element that combines with itself to generate the entire group are called \textbf{cyclic}. Recall that we can also have subgroups which are cyclic.
        \item<3-> Most groups we have worked with are not cyclic and have multiple generators. $\mathbb{Z}$, for example, has $\pm 1$ as a generator. $D_{4}$ has $R_{90}$ and $H$ as generators. 
        \item<4-> We call the set of generators of G the \textbf{generating set} which is usually called S. For $D_{4}$, $S = \{R_{90}, H\}$.
    \end{itemize}
\end{frame}

\begin{frame}
\frametitle{Generators and relations}
    \begin{itemize}
        \item<1-> Relation definition: the smallest set of equations that specify the conditions the generators must satisfy. In $D_{4}$ $R^{4} = H^{2} = (RH)^{2} = R_{0}$.
        \item<2-> Because this is the minimal set we can derive all additional relations from this set.
        \item<3-> If two groups are defined by the same relations, they are isomorphic.
        \item<4-> Describing a group by its generators and relations is called a \textbf{group presentation}. The presentation for $D_{4}$ is $\langle R, H | R^{4} = H^{2} = (RH)^{2} = R_{0}\rangle$.
        \item<4-> In order to define generators and relations rigorously, we need a beautiful concept called free groups. More information is available in the appendix if you are interested.
    \end{itemize}
\end{frame}

\begin{frame}
\frametitle{Cayley Digraphs of Groups}

\textit{Purpose:} Visualize a group given by generators and relations

These are directed graphs where vertices are group elements and arcs, which are the generators, connect the vertices. \medskip

\textit{Note:} Edges are typically used to describe vertex connections in \textbf{undirected} graphs while arcs describe connections in \textbf{directed graphs}.

\begin{figure}
    \centering
    \begin{subfigure}[b]{0.3\textwidth}
        \centering
        \includegraphics[width=\textwidth]{images/Z_6_DiGraph.png}
        \caption{$\mathbb{Z}_{6}$ DiGraph}
        \label{fig:Z_6 DiGraph}
    \end{subfigure}
    \hfill
    \begin{subfigure}[b]{0.3\textwidth}
        \centering
        \includegraphics[width=\textwidth]{images/A_4_DiGraph.png}
        \caption{$A_{4}$ DiGraph}
        \label{fig:A_4 DiGraph}
    \end{subfigure}
    \hfill
    \begin{subfigure}[b]{0.3\textwidth}
        \centering
        \includegraphics[width=\textwidth]{images/S_3.png}
        \caption{$S_{3}$ DiGraph}
        \label{fig:Z_6 DiGraph}
    \end{subfigure}
\end{figure}
\end{frame}

\begin{frame}
\frametitle{Cayley Digraphs of Groups}

Example 8 shows us that we can make DiGraphs for infinite groups too.


\begin{figure}
    \centering
    \begin{subfigure}[b]{0.3\textwidth}
        \centering
        \includegraphics[width=\textwidth]{images/Z_3_Z_2.png}
        \caption{$\mathbb{Z}_{3} \oplus \mathbb{Z}_{2}$ DiGraph}
        \label{fig:Z_3 x Z_2 DiGraph}
    \end{subfigure}
    \hfill
    \begin{subfigure}[b]{0.3\textwidth}
        \centering
        \includegraphics[width=\textwidth]{images/D_4_DiGraph.png}
        \caption{$D_{4}$ DiGraph}
        \label{fig:D_4 DiGraph}
    \end{subfigure}
    \hfill
    \begin{subfigure}[b]{0.3\textwidth}
        \centering
        \includegraphics[width=\textwidth]{images/Infinite_DiGraph.png}
        \caption{Infinite DiGraph}
        \label{fig: Infinite Vertices}
    \end{subfigure}
\end{figure}
\end{frame}

\begin{frame}
\frametitle{Cayley Digraphs of Groups}
\begin{block}{Definition: Cayley Digraph of a Group}
Let G be a finite group and let S be a set of generators for G. We define a digraph Cay(S:G), called the Cayley digraph of G with the generating set S, as follows:
    \begin{enumerate}
        \item Each element of G is a vertex of Cay(S:G)
        \item For each x and y in G, there is an arc from x to y if and only if xs = y for some s $\in$ S.
    \end{enumerate}
\end{block}
\end{frame}

\end{document}